\chapter{Avatar: Faithful Absence and CNAEXT Rendering}
\label{ch:avatar}

CNA has two independent avatar systems. The XNA-compatible classes reproduce the observable
behavior of a reference assembly whose Xbox Live content service is gone; that path is inert.
The \cnaclass{SkinnedModelEXT} path is a CNA extension that loads content and renders GPU-skinned
meshes. Calling the extension does not change the base API's state or semantics.

\begin{center}
\begin{tabularx}{\textwidth}{@{}lXX@{}}
\toprule
& \textbf{Base Avatar API} & \textbf{SkinnedModelEXT (CNAEXT)} \\
\midrule
Purpose & Preserve observed XNA surface without service content & Load and render CNA assets \\
Skeleton & Fixed 71-transform contract; zero transforms & Content-defined (bundled assets: 19 bones) \\
Draw & Validates, then performs no rendering & \cnaclass{GraphicsDevice}/\cnaclass{SkinnedEffect} draw \\
Activation & Normal API surface & Explicit opt-in \\
Evidence & Source translation and behavioral tests & Content, runtime, renderer, and pixel tests \\
\bottomrule
\end{tabularx}
\end{center}

``Faithful'' here is scoped to the translated behavior and tests. It does not mean that a C++
port is binary-identical to the managed assembly or that Xbox avatar services are available.

\section{Base API contract}

\cnaclass{AvatarRenderer} ignores both constructor arguments. Reading \texttt{State} sets and
returns \cnaclass{AvatarRendererState::Unavailable}; no path reaches \texttt{Ready} or
\texttt{Loading}. The \cnaclass{IAvatarAnimation} draw overload rejects null, reads its
transforms and expression, and delegates. The transform-array overload requires exactly 71
entries and then does nothing. \texttt{BindPose} throws \texttt{InvalidOperationException}
unless the unreachable \texttt{Ready} state is present.

\cnaclass{AvatarDescription} requires 1,021 bytes. CNA treats a correctly sized value as valid
when byte zero is nonzero; the other bytes remain opaque. The description getter copies the
bytes, height/body type default lazily, and the static \texttt{Changed} event is never raised.
Both \texttt{CreateRandom} overloads return the same all-zero invalid description. The body-type
overload validates its enum but does not use it when building the result.

\texttt{BeginGetFromGamer}/\texttt{EndGetFromGamer} complete synchronously and return that
invalid description. The caller owns and must delete the heap-allocated
\cnaclass{IAsyncResult}. \cnaclass{AvatarAnimation} supplies 71 zero transforms, zero duration,
and a default expression for all 31 presets. The preset is retained only as the CNAEXT clip name
returned by \texttt{GetRealClipNameEXT()}.

These results are expected when the proprietary service-supplied mesh, texture, animation, and
description formats are unavailable. FNA has no corresponding Avatar classes, so CNA's source
was translated from the XNA 4.0 reference assembly instead of inferred from FNA.

\section{SkinnedModelEXT contract}

\cnaclass{SkinnedModelEXT} is move-only because it owns GPU resources. It contains:

\begin{itemize}
  \item parent indices, local bind poses, and global inverse bind poses;
  \item named clips with per-bone translation, scale, and quaternion-rotation tracks;
  \item skinned mesh parts and their vertex, index, material, and texture resources.
\end{itemize}

Its skeleton is unrelated to \cnaclass{AvatarRenderer}'s 71-entry validation array. The
bundled male and female assets use 19 bones; other assets may use different counts up to
\cnaclass{SkinnedEffect}'s 72-bone draw limit.

\texttt{AvatarRenderer::EnableRealRenderingEXT(device, model)} creates the extension's
\cnaclass{SkinnedEffect}. \texttt{DrawRealEXT(clip, position, loop)} samples a clip and draws
the model, or throws \texttt{InvalidOperationException} if opt-in has not occurred. It does not
make \texttt{State} ready or give the base \texttt{Draw()} a renderer. Conventional roots are
\texttt{avatar/male/avatar} and \texttt{avatar/female/avatar}.

\begin{lstlisting}
AvatarDescription desc = AvatarDescription::CreateRandom();
AvatarRenderer renderer(&desc); // base state remains Unavailable

auto& device = getGraphicsDeviceProperty();
auto model = getContentProperty().Load<std::shared_ptr<SkinnedModelEXT>>(
    "avatar/male/avatar");
renderer.EnableRealRenderingEXT(device, model);
renderer.SetAppearanceEXT(AvatarAppearanceEXT());

renderer.setWorldProperty(Matrix::getIdentityProperty());
renderer.setViewProperty(camera.View);
renderer.setProjectionProperty(camera.Projection);
renderer.DrawRealEXT("Stand0", elapsedClipTime, true);
\end{lstlisting}

The fragment is intended for a \cnaclass{Game} subclass with the Avatar, content, graphics, and
\texttt{<memory>} headers already included.

\section{Animation and palette computation}

\texttt{ComputeBoneTransformsEXT} first initializes every local transform from the bind pose,
then samples matching tracks. Translation and scale use linear interpolation; rotation uses
\texttt{Quaternion::Slerp}. It composes children with already computed parents in stored
topological order and finally multiplies each global transform by its global inverse bind pose.
The result is the palette consumed by \cnaclass{SkinnedEffect}.

\begin{lstlisting}
SkinnedModelEXT model;
model.BoneCount = 2;
model.ParentBoneIndices = {-1, 0};
model.BindPoseLocal = {
    Matrix::getIdentityProperty(),
    Matrix::CreateTranslation(Vector3(0, 1, 0))};
model.InverseBindPoseGlobal = {
    Matrix::getIdentityProperty(), Matrix::getIdentityProperty()};

BoneTrackEXT track;
track.BoneIndex = 0;
track.Keys.push_back(KeyframeEXT{
    TimeSpan::FromSeconds(0), Vector3(0, 0, 0)});
track.Keys.push_back(KeyframeEXT{
    TimeSpan::FromSeconds(1), Vector3(2, 0, 0)});

AnimationClipEXT clip;
clip.Duration = TimeSpan::FromSeconds(1);
clip.Tracks.push_back(track);
model.Clips["Move"] = clip;

std::vector<Matrix> palette;
model.ComputeBoneTransformsEXT(
    "Move", TimeSpan::FromSeconds(0.5), false, palette);
// palette[0].Translation.X == 1.0f
\end{lstlisting}

Current guards address four earlier failures:

\begin{itemize}
  \item looped time uses tick modulo instead of repeated subtraction, including negative time;
  \item each non-root parent must precede its child, which also rejects cycles;
  \item skeleton arrays must match \texttt{BoneCount} before indexing;
  \item empty or out-of-range animation tracks are skipped without touching memory.
\end{itemize}

The time regression uses a position 1,000,000,000.5 clip durations from the origin and a
100-ms bound, specifically detecting the former duration-by-duration loop.

\section{Lighting, appearance, and wardrobe}

The extension uses one directional light plus ambient light. It disables the other two stock
directional lights and zeroes specular and emissive terms before applying the caller's values.
Default ambient and key-light values keep an avatar visible without configuration.

\cnaclass{AvatarAppearanceEXT} is a CNA tint model, not a reconstruction of the proprietary
appearance format. Case-sensitive part-name substrings select \texttt{Hair}, \texttt{Shirt},
\texttt{Pants}, or \texttt{Shoes}; unmatched names receive skin tint. It changes colors, not
geometry or per-part textures.

\texttt{AttachPartEXT} requires a matching bone count. Reattaching a part with the same name
first calls \texttt{RemovePartEXT}, which releases the old owned vertex/index buffers, mesh
part, and texture before inserting the replacement. This closes both the former overlapping-part
behavior and its GPU-resource leak.

\begin{historynote}{Lighting defects found through avatar rendering}
An early draw path applied ambient light before \texttt{EnableDefaultLighting()}, which then
overwrote it. A separate EasyGL shader defect multiplied emissive color by diffuse color twice,
turning the ambient term into $ambient\times diffuse^2$. That defect affected all related
skinned, vertex-lit, and environment-mapped variants. A pixel test's near-black fraction fell
from 4.1--6.0\% to 0\% after the shader formula was aligned with FNA's lighting expression.
\end{historynote}

\section{Content pipeline}

The \texttt{.skinnedmodel.json}, \texttt{.skeleton.bin}, and \texttt{.clip.bin} pipeline loads
the extension assets. Three corrected implementation traps are relevant to other importers:

\begin{itemize}
  \item manifest-relative paths must resolve from the manifest directory, not the content root;
  \item chained \texttt{reader.Read<float>()} constructor arguments relied on unspecified C++
        argument evaluation order, so components are now read into named locals;
  \item hierarchy reordering must remap inverse bind matrices and per-vertex joint indices as
        well as parent indices.
\end{itemize}

Chapter~\ref{ch:skinning-animation} covers the corresponding glTF joint-space rules.

\section{Open garment interpenetration}

The procedural avatar still shows shirt and skin fragments around the raised shoulder and
chest in the \texttt{Wave} clip. Measurement found the same 64 of 228 shirt vertices crossing
the body in both neutral and Wave poses; the animation changes which crossings face the camera.
This locates the remaining defect in the rest-pose geometry, not clip sampling.

The largest visible intersections are the collar and cuffs. The shirt's \texttt{Spine1} shell
has radius 0.176 and spans $z=1.40$ to $1.576$; the head sphere has radius 0.15 and center
$z=1.60$, placing its bottom at $z=1.45$. Their calculated overlap of about 0.126 m agrees with
the measured 0.131 m collar penetration. A waist overlap is mutually hidden by garments.

Four attempted remedies did not remove the cause:

\begin{center}
\begin{tabularx}{\textwidth}{@{}lX@{}}
\toprule
\textbf{Attempt} & \textbf{Observed result} \\
\midrule
Narrow garment blend radius & No clear visual improvement \\
Trim capsule boundary caps & Pants crossings rose from 21 to 39 \\
Flat cylinders plus joint spheres & More speckle, 28 new groin crossings, lower brightness \\
32 instead of 12 segments & Smoother edges, same fragments, 4.2$\times$ vertex data \\
\bottomrule
\end{tabularx}
\end{center}

The open choices are to size collar and cuff shells against the body radii along their full
span, accept the limitation of convex per-bone placeholder shells, or explicitly accept the
tessellation cost. The extension is therefore renderer-functional but not visually defect-free.

\section{Renderer and evidence status}

EasyGL and Vulkan run registered avatar integration programs with pixel readback. Their
synthetic one-bone quads isolate palette/effect/draw wiring; they do not validate the entire
asset converter. Demo programs add full humanoid content, clip changes, appearance, attachment,
wardrobe replacement, API-boundary comparison, and two-process network synchronization.

BGFX contains relevant pipeline code but has no registered avatar integration result at the
pinned revision; existing harnesses stop earlier at unimplemented 3D depth/blend setup.
\texttt{SDL\_RENDERER} rejects the required 3D resources. No conclusion for either renderer
should be inferred from EasyGL or Vulkan pixels. Appendix~\ref{ch:verification-tiers} defines
the evidence terms used here.
